% ─────────────────────────────────────────────────────────────────────────────
\section{Evolution Package (\texttt{src/evolution/})}
\label{sec:evolution}
% ─────────────────────────────────────────────────────────────────────────────

The evolution package implements the NEAT algorithm: genome representation, neural
network construction and inference, speciation, mutation, crossover, and the generation
evaluation loop.

\begin{figure}[H]
  \centering
  \includegraphics[width=\textwidth]{evolution_classes}
  \caption{Evolution package class diagram.}
  \label{fig:evo}
\end{figure}

% ─────────────────────────────────────────────────────────────────────────────
\subsection{\texttt{NodeType} (enum class)}

\begin{lstlisting}
enum class NodeType { Input, Hidden, Output, Bias };
\end{lstlisting}

\texttt{Bias} nodes have a fixed activation of 1.0 and appear in every initial genome.

% ─────────────────────────────────────────────────────────────────────────────
\subsection{\texttt{NodeGene} and \texttt{ConnectionGene} (structs)}

\begin{table}[H]
\centering\small
\caption{Gene struct fields}
\begin{tabularx}{\textwidth}{@{}lllX@{}}
\toprule
\textbf{Struct} & \textbf{Field} & \textbf{Type} & \textbf{Description} \\
\midrule
\texttt{NodeGene}       & \texttt{id}         & \texttt{uint32\_t} & Global node identifier \\
                        & \texttt{type}       & \texttt{NodeType}  & Input/Hidden/Output/Bias \\
\midrule
\texttt{ConnectionGene} & \texttt{in\_node}   & \texttt{uint32\_t} & Source node id \\
                        & \texttt{out\_node}  & \texttt{uint32\_t} & Target node id \\
                        & \texttt{weight}     & \texttt{float}     & Synaptic weight \\
                        & \texttt{enabled}    & \texttt{bool}      & False = structurally present but not expressed \\
                        & \texttt{innovation} & \texttt{uint32\_t} & Historical marking for alignment during crossover \\
\bottomrule
\end{tabularx}
\end{table}

% ─────────────────────────────────────────────────────────────────────────────
\subsection{\texttt{Genome} (class)}

Encodes a variable-topology neural network as gene lists.

\begin{table}[H]
\centering\small
\caption{\texttt{Genome} interface}
\begin{tabularx}{\textwidth}{@{}llX@{}}
\toprule
\textbf{Method} & \textbf{Return} & \textbf{Description} \\
\midrule
\texttt{Genome(num\_inputs, num\_outputs)} & — & Create minimal genome (inputs + bias + outputs; no hidden) \\
\texttt{add\_node(node)} & \texttt{void} & Append \texttt{NodeGene} \\
\texttt{add\_connection(conn)} & \texttt{void} & Append \texttt{ConnectionGene} \\
\texttt{has\_connection(from, to)} & \texttt{bool} & Check structural connectivity \\
\texttt{has\_node(id)} & \texttt{bool} & Membership test \\
\texttt{max\_node\_id()} & \texttt{uint32\_t} & Largest node id in genome \\
\texttt{nodes()} & \texttt{const/mutable ref} & Gene list accessors \\
\texttt{connections()} & \texttt{const/mutable ref} & Gene list accessors \\
\texttt{fitness() / set\_fitness(f)} & \texttt{float} & Raw fitness from evaluation \\
\texttt{adjusted\_fitness() / set\_adjusted\_fitness(f)} & \texttt{float} & Shared fitness (divided by species size) \\
\texttt{num\_inputs(), num\_outputs()} & \texttt{int} & Topology constants \\
\texttt{complexity()} & \texttt{int} & \texttt{nodes\_.size() + enabled\_connection\_count} \\
\texttt{compatibility\_distance(a,b,c1,c2,c3)} & \texttt{float} & Static; compatibility metric (see below) \\
\texttt{to\_json()} & \texttt{string} & Serialize to JSON string \\
\texttt{from\_json(str)} & \texttt{Genome} & Static factory from JSON string \\
\bottomrule
\end{tabularx}
\end{table}

\subsubsection{Compatibility Distance Algorithm}

Defined by Stanley \& Miikkulainen~\cite{stanley2002}:
\[
  \delta = \frac{c_1 \cdot E + c_2 \cdot D}{\max(N_1, N_2)} + c_3 \cdot \bar{W}
\]
where $E$ = excess gene count, $D$ = disjoint gene count, $\bar{W}$ = mean weight
difference of matching genes, and $N_i$ = genome size.  Genes are aligned by innovation
number.

\subsubsection{JSON Serialization Format}

\begin{lstlisting}[language={},caption={Genome JSON schema}]
{
  "nodes": [{"id": 0, "type": "input"}, ...],
  "connections": [
    {"in": 0, "out": 16, "weight": 0.42, "enabled": true, "innovation": 3}, ...
  ]
}
\end{lstlisting}

\textbf{Private members:} \texttt{nodes\_}, \texttt{connections\_},
\texttt{num\_inputs\_}, \texttt{num\_outputs\_}, \texttt{fitness\_},
\texttt{adjusted\_fitness\_}.

% ─────────────────────────────────────────────────────────────────────────────
\subsection{\texttt{NeuralNetwork} (class)}

Builds a runnable feedforward network from a \texttt{Genome}.

\subsubsection{Construction Algorithm}
\begin{enumerate}[nosep]
  \item Build \texttt{nodes\_} list from enabled nodes in genome.
  \item Topological sort: Kahn's algorithm on the enabled connection DAG.  DFS cycle
    detection is applied to any new hidden connections before insertion.  Result stored
    in \texttt{evaluation\_order\_} (list of \texttt{uint32\_t} node ids in activation
    order).
  \item Precompute \texttt{incoming\_[i]}: for each node at eval-order index $i$,
    collect all \texttt{(source\_idx, weight)} pairs from enabled connections.
  \item Allocate \texttt{values\_} float vector of length \texttt{nodes\_.size()}.
\end{enumerate}

\subsubsection{\texttt{activate(inputs)} Algorithm}
\begin{enumerate}[nosep]
  \item Copy \texttt{inputs} into \texttt{values\_} at input node indices.
  \item Set bias node value to 1.0.
  \item Iterate \texttt{evaluation\_order\_}:
    \[
      v_i = f\!\left(\sum_{(j,w) \in \mathit{incoming}[i]} w \cdot v_j\right)
    \]
    where $f$ is the configured activation function.
  \item Return output node values in order.
\end{enumerate}

\subsubsection{Activation Functions}

\begin{lstlisting}
enum class ActivationFn { Sigmoid, Tanh, ReLU };
\end{lstlisting}

\begin{table}[H]
\centering\small
\caption{Activation function formulas}
\begin{tabular}{@{}lll@{}}
\toprule
\textbf{Name} & \textbf{Formula} & \textbf{Output range} \\
\midrule
\texttt{Sigmoid} & $1/(1+e^{-x})$ & $(0,1)$ \\
\texttt{Tanh}    & $\tanh(x)$     & $(-1,1)$ \\
\texttt{ReLU}    & $\max(0,x)$    & $[0,\infty)$ \\
\bottomrule
\end{tabular}
\end{table}

\begin{table}[H]
\centering\small
\caption{\texttt{NeuralNetwork} interface}
\begin{tabularx}{\textwidth}{@{}llX@{}}
\toprule
\textbf{Method} & \textbf{Return} & \textbf{Description} \\
\midrule
\texttt{NeuralNetwork(genome, activation\_fn)} & — & Construct; run build algorithm \\
\texttt{activate(inputs)} & \texttt{vector<float>} & Forward pass; returns output activations \\
\texttt{num\_nodes(), num\_connections()} & \texttt{int const} & Topology counts \\
\texttt{num\_inputs(), num\_outputs()} & \texttt{int const} & I/O sizes \\
\texttt{raw\_nodes()} & \texttt{const vector<Node>\&} & All node descriptors \\
\texttt{eval\_order()} & \texttt{const vector<uint32\_t>\&} & Activation sequence (for GPU packing) \\
\texttt{incoming()} & \texttt{const vector<vector<pair<int,float>>>\&} & Precomputed adjacency \\
\texttt{node\_index\_map()} & \texttt{const unordered\_map<uint32\_t,int>\&} & node\_id → value index \\
\texttt{last\_activations()} & \texttt{const vector<float>\&} & Activations from most recent \texttt{activate()} call \\
\bottomrule
\end{tabularx}
\end{table}

% ─────────────────────────────────────────────────────────────────────────────
\subsection{\texttt{Species} (class)}

Groups genomes by structural similarity.  Manages stagnation detection and adjusted
fitness for fair reproduction allocation.

\begin{table}[H]
\centering\small
\caption{\texttt{Species} private members}
\begin{tabularx}{\textwidth}{@{}llX@{}}
\toprule
\textbf{Member} & \textbf{Type} & \textbf{Description} \\
\midrule
\texttt{id\_}                           & \texttt{int}            & Unique species id \\
\texttt{representative\_}               & \texttt{Genome}         & Reference genome for compatibility tests \\
\texttt{members\_}                      & \texttt{vector<Genome*>}& Non-owning pointers into \texttt{EvolutionManager::population\_} \\
\texttt{average\_fitness\_}             & \texttt{float}          & Mean raw fitness of members \\
\texttt{total\_adjusted\_fitness\_}     & \texttt{float}          & Sum of adjusted fitnesses \\
\texttt{best\_fitness\_ever\_}          & \texttt{float}          & Highest fitness seen since species creation \\
\texttt{generations\_without\_improvement\_} & \texttt{int}       & Incremented each generation max fitness does not improve \\
\texttt{fitness\_history\_}             & \texttt{deque<float>}   & Rolling history of average fitness \\
\bottomrule
\end{tabularx}
\end{table}

\subsubsection{Key Algorithms}

\textbf{\texttt{adjust\_fitness()}}: Divides each member's raw fitness by species size
(explicit fitness sharing):
\[
  f'_i = f_i / |\mathit{members}|
\]

\textbf{\texttt{is\_stagnant(limit)}}: Returns
\texttt{generations\_without\_improvement\_ >= limit}.

\textbf{\texttt{update\_best\_fitness()}}: If the current maximum member fitness exceeds
\texttt{best\_fitness\_ever\_}, update it and reset the stagnation counter; otherwise
increment it.

\textbf{\texttt{update\_representative()}}: Set \texttt{representative\_} to the
highest-fitness member (used for next-generation compatibility tests).

% ─────────────────────────────────────────────────────────────────────────────
\subsection{\texttt{InnovationTracker} (class)}

Ensures historical markings are consistent within a generation, enabling correct NEAT
crossover alignment.

\begin{table}[H]
\centering\small
\caption{\texttt{InnovationTracker} interface}
\begin{tabularx}{\textwidth}{@{}llX@{}}
\toprule
\textbf{Method} & \textbf{Return} & \textbf{Description} \\
\midrule
\texttt{init\_from\_population(population)} & \texttt{void} & Scan all genomes; set counters to max innovation / node ids found \\
\texttt{get\_innovation(in, out)} & \texttt{uint32\_t} & Look up or create innovation for edge (in→out) within this generation \\
\texttt{next\_node\_id()} & \texttt{uint32\_t} & Return and increment the global node counter \\
\texttt{reset\_generation()} & \texttt{void} & Clear \texttt{generation\_innovations\_} (called at start of \texttt{evolve()}) \\
\texttt{innovation\_count(), node\_count()} & \texttt{uint32\_t const} & Current counter values \\
\texttt{set\_counters(innov, node)} & \texttt{void} & Restore from checkpoint \\
\bottomrule
\end{tabularx}
\end{table}

\textbf{Private:} \texttt{generation\_innovations\_: map<pair<uint32\_t,uint32\_t>,uint32\_t>}
— the within-generation map that guarantees the same structural mutation receives the
same innovation number across all genomes in one generation.

% ─────────────────────────────────────────────────────────────────────────────
\subsection{\texttt{Mutation} (static class)}

All methods are \texttt{static}.

\begin{table}[H]
\centering\small
\caption{\texttt{Mutation} operator algorithms}
\begin{tabularx}{\textwidth}{@{}lX@{}}
\toprule
\textbf{Method} & \textbf{Algorithm} \\
\midrule
\texttt{mutate\_weights(genome, rng, power)} &
  For each connection: with probability 0.9 apply Gaussian perturbation
  $w \leftarrow w + \mathcal{N}(0, \mathit{power})$; with probability 0.1 replace
  with uniform sample from $[-2, 2]$. \\
\texttt{add\_connection(genome, rng, tracker)} &
  Sample random unconnected node pair; check no cycle (DFS); assign innovation via
  \texttt{tracker.get\_innovation()}; initialize weight $\sim \mathcal{U}(-1,1)$. \\
\texttt{add\_node(genome, rng, tracker, max\_hidden)} &
  Select random enabled connection; disable it; insert new hidden node
  (\texttt{tracker.next\_node\_id()}); add connection \texttt{in → new} with weight 1.0;
  add connection \texttt{new → out} with old weight. \\
\texttt{toggle\_connection(genome, rng)} &
  Pick random connection; flip \texttt{enabled} flag. \\
\texttt{mutate(genome, rng, config, tracker)} &
  Apply each operator independently: \texttt{mutate\_weights} with
  \texttt{mutation\_rate}, \texttt{add\_node} with \texttt{add\_node\_rate},
  \texttt{add\_connection} with \texttt{add\_connection\_rate},
  \texttt{toggle\_connection} at a fixed low rate. \\
\bottomrule
\end{tabularx}
\end{table}

% ─────────────────────────────────────────────────────────────────────────────
\subsection{\texttt{Crossover} (static class)}

\begin{lstlisting}
static Genome crossover(const Genome& parent_a,
                        const Genome& parent_b, Random& rng);
\end{lstlisting}

\textbf{Algorithm:} Align genes by innovation number.
\begin{enumerate}[nosep]
  \item \textbf{Matching genes:} Pick from either parent with probability 0.5.
  \item \textbf{Disjoint / excess genes:} Always inherit from the fitter parent.
  \item \textbf{Disabled genes:} If disabled in either contributing parent, disable
    in child with probability 0.75.
  \item Inherit \texttt{num\_inputs} and \texttt{num\_outputs} from the fitter parent.
\end{enumerate}

% ─────────────────────────────────────────────────────────────────────────────
\subsection{\texttt{EvolutionManager} (class)}

Top-level controller for NEAT: holds the population, networks, species, and orchestrates
the evaluate → compute\_fitness → evolve cycle.

\begin{table}[H]
\centering\small
\caption{\texttt{EvolutionManager} private members}
\begin{tabularx}{\textwidth}{@{}llX@{}}
\toprule
\textbf{Member} & \textbf{Type} & \textbf{Description} \\
\midrule
\texttt{config\_}     & \texttt{SimulationConfig}                    & Run parameters \\
\texttt{rng\_}        & \texttt{Random\&}                            & Shared RNG (owned by main) \\
\texttt{tracker\_}    & \texttt{InnovationTracker}                   & Innovation numbering \\
\texttt{population\_} & \texttt{vector<Genome>}                      & Full population by value \\
\texttt{species\_}    & \texttt{vector<Species>}                     & Current species list \\
\texttt{networks\_}   & \texttt{vector<unique\_ptr<NeuralNetwork>>}  & One NN per genome \\
\texttt{generation\_} & \texttt{int}                                 & Current generation index \\
\texttt{num\_inputs\_, num\_outputs\_} & \texttt{int}             & Determined at \texttt{initialize()} \\
\texttt{use\_gpu\_}   & \texttt{bool}                                & GPU inference enabled \\
\texttt{gpu\_batch\_} & \texttt{unique\_ptr<GpuBatch>}               & Device memory (CUDA build only) \\
\texttt{tick\_callback\_} & \texttt{TickCallback}                    & Optional per-tick hook \\
\bottomrule
\end{tabularx}
\end{table}

\begin{table}[H]
\centering\small
\caption{\texttt{EvolutionManager} public interface}
\begin{tabularx}{\textwidth}{@{}llX@{}}
\toprule
\textbf{Method} & \textbf{Return} & \textbf{Description} \\
\midrule
\texttt{EvolutionManager(config, rng)} & — & Store config and RNG reference \\
\texttt{initialize(num\_inputs, num\_outputs)} & \texttt{void} & Create minimal genomes; build networks \\
\texttt{evaluate\_generation(sim)} & \texttt{void} & Run one generation (see below) \\
\texttt{evolve()} & \texttt{void} & Speciate → remove stagnant → reproduce \\
\texttt{compute\_fitness(sim)} & \texttt{void} & Write fitness into each genome (see below) \\
\texttt{population()} & \texttt{const/mutable ref} & Access genome vector \\
\texttt{generation()} & \texttt{int const} & Current generation number \\
\texttt{species()} & \texttt{const vector<Species>\&} & Current species list \\
\texttt{networks()} & \texttt{const vector<unique\_ptr<NN>>\&} & Built network list \\
\texttt{set\_tick\_callback(cb)} & \texttt{void} & Install \texttt{TickCallback} \\
\texttt{clear\_tick\_callback()} & \texttt{void} & Remove callback \\
\texttt{enable\_gpu(flag)} & \texttt{void} & Toggle GPU inference path \\
\texttt{gpu\_enabled()} & \texttt{bool const} & Query current GPU state \\
\texttt{save\_checkpoint(path, rng)} & \texttt{void} & Serialize to JSON (see below) \\
\texttt{load\_checkpoint(path, rng)} & \texttt{bool} & Restore from checkpoint; false if not found \\
\bottomrule
\end{tabularx}
\end{table}

\subsubsection{\texttt{evaluate\_generation()} Algorithm}

\begin{enumerate}[nosep]
  \item Call \texttt{sim.reset()} and \texttt{sim.initialize()}.
  \item For each tick $t \in [0, \mathit{generation\_ticks})$:
    \begin{enumerate}[nosep,label=\alph*.]
      \item For each alive agent $i$: call \texttt{sim.get\_sensors(i)};
        run \texttt{networks\_[i]->activate()}; decode output to direction vector;
        call \texttt{sim.apply\_action(i, direction, dt)}.
      \item Call \texttt{sim.tick(dt)}.
      \item If \texttt{tick\_callback\_} is set, invoke it with $(t, \mathit{sim})$.
    \end{enumerate}
  \item On GPU path: \texttt{gpu\_batch\_->pack\_inputs()} before agent loop;
    \texttt{batch\_neural\_inference()} launches CUDA kernel;
    \texttt{gpu\_batch\_->unpack\_outputs()} retrieves results.
\end{enumerate}

\subsubsection{\texttt{compute\_fitness()} Algorithm}

For each genome $i$ paired with \texttt{agents\_[i]}:
\[
  f_i = w_s \cdot \frac{\mathit{age}}{T} + w_k \cdot \mathit{kills} + w_e \cdot \frac{\mathit{energy}}{E_0}
        + w_d \cdot \frac{\mathit{dist}}{d_{\max}} - w_c \cdot \mathit{complexity}
\]
where $T$ = \texttt{generation\_ticks}, $E_0$ = \texttt{initial\_energy},
$d_{\max} = \sqrt{w^2+h^2}$, and the $w_*$ constants come from
\texttt{SimulationConfig}.

\subsubsection{\texttt{evolve()} Algorithm}
\begin{enumerate}[nosep]
  \item \texttt{speciate()}: For each genome, find the first compatible species
    (representative comparison).  No match → new species.  Update
    \texttt{update\_best\_fitness()} on each species.
  \item \texttt{remove\_stagnant\_species()}: Remove species where
    \texttt{is\_stagnant(stagnation\_limit)} is true (except the top-fitness species).
  \item Compute offspring allocation via largest-remainder method proportional to
    total adjusted fitness.
  \item \texttt{reproduce()}: For each species, elitism (top genome survives);
    remainder via \texttt{Crossover::crossover} + \texttt{Mutation::mutate} using
    tournament selection within species.
  \item \texttt{tracker\_.reset\_generation()}.
  \item \texttt{build\_networks()}.
  \item Increment \texttt{generation\_}.
\end{enumerate}

\subsubsection{Checkpoint Format}

JSON file with fields: \texttt{generation}, \texttt{population} (array of genome JSON),
\texttt{rng\_state} (hex string of engine state), \texttt{species\_stagnation} (array of
\texttt{\{id, best\_ever, gens\_no\_improvement\}}), \texttt{innovation\_counter},
\texttt{node\_counter}.

\subsubsection{\texttt{TickCallback} Type Alias}

\begin{lstlisting}
using TickCallback =
    std::function<void(int tick, const SimulationManager& sim)>;
\end{lstlisting}

\texttt{main.cpp} uses this to wire \texttt{Logger::log\_tick}, live visualization
updates, and per-generation benchmarking probes into the evaluation loop without creating
any dependency from \texttt{evolution/} on \texttt{data/} or \texttt{visualization/}.
